% Figure: oblique two-vehicle collision momentum vector diagram. Car A
% travels east, car B travels north; they lock and move off at a common
% heading. The vector triangle shows total momentum conserved.
\begin{figure}[htbp]
\centering
\begin{tikzpicture}[
  pvec/.style={draw=engblue, -{Stealth}, very thick},
  rvec/.style={draw=engred, -{Stealth}, very thick},
  lbl/.style={font=\scriptsize},
]
% scale: 1 unit = 10000 kg m/s.
% Car A: m=1400, v=14 east  -> pA = 19600 -> 1.96 units east
% Car B: m=1000, v=18 north -> pB = 18000 -> 1.80 units north
% total p = (1.96, 1.80), magnitude 2.66 units = 26600 kg m/s
% origin
\coordinate (O) at (0,0);
% pA along x
\draw[pvec] (O) -- (3.92,0) node[midway, below, lbl] {$p_A=19\,600$};
% pB along y from tip of pA (vector triangle)
\draw[pvec] (3.92,0) -- (3.92,3.60) node[midway, right, lbl] {$p_B=18\,000$};
% resultant from O to tip
\draw[rvec] (O) -- (3.92,3.60) node[midway, above left, lbl, engred]
  {$P=26\,600$};
% angle annotation
\node[lbl, anchor=west] at (1.4,1.0) {$\phi=42.6^{\circ}$};
% small axis compass
\draw[->, enggray] (-0.6,-0.6) -- (-0.6,0.2) node[above, lbl, enggray] {N};
\draw[->, enggray] (-0.6,-0.6) -- (0.2,-0.6) node[right, lbl, enggray] {E};
% labels for the two incoming cars (schematic)
\node[lbl, engblue, anchor=north] at (1.96,-0.7) {car A: east};
\node[lbl, engblue, anchor=west] at (4.05,1.8) {car B: north};
% post-impact wreck heading note
\node[lbl, engred, anchor=west] at (2.3,3.0) {wreck heading};
\end{tikzpicture}
\caption[Oblique two-vehicle collision: momentum vector triangle]{The
momentum vector triangle for the oblique collision in the
accident-reconstruction case study. Car A ($1{,}400\,\si{\kilogram}$)
travels east at $14\,\si{\meter\per\second}$, contributing
$p_A=19{,}600\,\si{\kilogram\meter\per\second}$; car B
($1{,}000\,\si{\kilogram}$) travels north at $18\,\si{\meter\per
\second}$, contributing $p_B=18{,}000\,\si{\kilogram\meter\per
\second}$. The two perpendicular momenta add head to tail, so the
locked wreck departs along the resultant $P=\sqrt{p_A^{2}+p_B^{2}}=
26{,}600\,\si{\kilogram\meter\per\second}$ at heading
$\phi=\arctan(p_B/p_A)=42.6^{\circ}$ north of east. The post-impact
skid direction recorded at the scene is the physical check on this
heading, and a mismatch is the first sign that a stated incoming speed
is wrong.}
\label{fig:vol03ch04:oblique-collision}
\end{figure}
